\section{Methodology}
\label{sec:method}

\subsection{Setup}
\label{sec:setup}

Every experiment targets the VPR architecture \code{k6\_frac\_N10\_frac\_chain\_mem32K\_40nm} ($K=6$, ten BLEs per CLB, carry chains, 32\,Kb memories, multipliers). The netlist is the VTR-dialect BLIF that Odin~III's front end writes: \code{read\_verilog --nolatches --vtr-compat}, \code{proc}, \code{fsm}, \code{opt}, \code{flatten}, \code{opt -full}, \code{wreduce}, \code{hardmap}, \code{lower}, \code{abc --lut 6}, \code{stats}, \code{write\_blif --vtr --lut-k 6 --name-style provenance}. The seeding baseline runs on VPR \code{v8.0.0+18416} with one deterministic VPR seed per configuration. The packer experiments use VTR commit \code{3c9a4d23b2} (VTR~\cite{rose2012vtr}, VTR~9~\cite{elgammal2025vtr9}) with three VPR seeds per configuration unless the experiment says otherwise. The designs are: the 15 VTR~8.0 designs of Table~\ref{tab:benchmarks}; the generated yardstick designs (72, 105 and 84 designs, Section~\ref{sec:i-yardstick}); and for the option sweep 35 designs, the 20 MCNC circuits and the 15 designs of Table~\ref{tab:benchmarks}. Every number in the paper traces to a committed file, named in a comment at the number in the source, in the research log (\code{docs/paper/hypotheses.md}) or in \code{data/}. Raw run directories are not committed; their summaries are.

\subsection{Benchmarks}
\label{sec:benchmarks}

Table~\ref{tab:benchmarks} reports each design's construction and its post-technology-mapping size. The module count is the front end's, while the CLB count and the grid are VPR's own numbers for the packed netlist.

% BEGIN tab:benchmarks prose
The 15 designs span two orders of magnitude after technology mapping:
from 13 CLBs (ch\_intrinsics) to 1956 (stereovision2),
4843 CLBs in total. The mean is 322.9 CLBs and the median 125;
the mean sits well above the median because three stereovision designs carry most of the logic.
VPR auto-sizes a square grid to the packed netlist, with a mean side of 27.2 tiles
(median 20) and a mean of 1076.4 grid tiles (median
400). Grid size follows the largest resource demand, not the CLB count:
I/O-heavy and memory-heavy designs such as mkDelayWorker32B get a large die for few CLBs.
% END tab:benchmarks prose

\begin{table*}[t]
\caption{Benchmark construction and post-technology-mapping size. All fifteen
designs are from the VTR~8.0 benchmark set
(\code{vtr\_flow/benchmarks/verilog}). Modules is the front end's module
count, CLBs the \code{clb} blocks VPR packs, Memory and Mult the memory and
\code{mult\_36} blocks, and Grid the auto-sized device grid with its tile
count in parentheses. Sizes are read from each design's baseline run log and
kept in \code{data/benchmark-sizes.json}.}
\label{tab:benchmarks}
\centering
\footnotesize
\begin{tabular}{@{}l l r r r r c@{}}
\toprule
Design & Suite & Modules & CLBs & Memory & Mult & Grid (tiles) \\
\midrule
ch\_intrinsics & VTR 8.0 & 2 & 13 & 1 & 0 & $10\times$10 (100) \\
stereovision3 & VTR 8.0 & 1 & 15 & 0 & 0 & $7\times$7 (49) \\
diffeq2 & VTR 8.0 & 1 & 22 & 0 & 5 & $16\times$16 (256) \\
diffeq1 & VTR 8.0 & 1 & 37 & 0 & 5 & $16\times$16 (256) \\
mkPktMerge & VTR 8.0 & 7 & 21 & 15 & 0 & $28\times$28 (784) \\
raygentop & VTR 8.0 & 15 & 143 & 1 & 6 & $19\times$19 (361) \\
sha & VTR 8.0 & 1 & 142 & 0 & 0 & $17\times$17 (289) \\
boundtop & VTR 8.0 & 13 & 54 & 1 & 0 & $15\times$15 (225) \\
mkSMAdapter4B & VTR 8.0 & 8 & 125 & 5 & 0 & $20\times$20 (400) \\
or1200 & VTR 8.0 & 16 & 252 & 2 & 1 & $27\times$27 (729) \\
stereovision0 & VTR 8.0 & 25 & 707 & 0 & 0 & $33\times$33 (1089) \\
stereovision1 & VTR 8.0 & 16 & 788 & 0 & 43 & $40\times$40 (1600) \\
mkDelayWorker32B & VTR 8.0 & 16 & 74 & 43 & 0 & $50\times$50 (2500) \\
stereovision2 & VTR 8.0 & 24 & 1956 & 0 & 192 & $82\times$82 (6724) \\
blob\_merge & VTR 8.0 & 2 & 494 & 0 & 0 & $28\times$28 (784) \\
\bottomrule
\end{tabular}
\end{table*}


\subsection{Metrics}
\label{sec:metrics}

\pp{Area (low is better).} The LUT count of the Odin~III netlist, the number of logic CLBs after packing, the die grid VPR chooses, and the minimum routable channel width.

\pp{Speed (high is better): critical and near-critical.} The speed numbers are the critical-path delay (CPD) in nanoseconds and the near-critical statistics below. The routed wirelength (WL) is reported beside them as a proxy for wire delay and routing area; it is the primary metric of the packer experiments because their plans named it. A design is not fast only because its single slowest path is fast, since the next slowest path becomes the CPD once the first is sped up. We therefore also look at near-critical connections, those within 10\% of the CPD in the routed timing report. VPR defines the criticality of a connection from its slack, normalised by the CPD~\cite{luu2014thesis}; VPR~9 computes it relaxed per timing constraint~\cite{wainberg2015relaxed}. Only one near-critical statistic exists in this paper (the experiment of Section~\ref{sec:e-census}); worst slack per object is not yet measured (Section~\ref{sec:mimir}).

\pp{Power.} Power is not measured in any experiment of this paper. We make no claim about it, and we cite no relation of power to area.

\subsection{Hypotheses, experiments and verdicts}
\label{sec:verdicts}

Each experiment subsection of Section~\ref{sec:experiments} has the same four parts: the hypothesis in one sentence, the experiment (designs, arms, seeds), the result (numbers and the test), and the verdict. The verdict rule of each experiment is a numeric rule written in its plan file under \code{docs/plan/}; a rule is only changed by a dated entry that keeps the old text, and where a rule was changed the original verdict is the one reported. The four verdicts are:

\begin{itemize}
\item \textsc{useful}: supported under the experiment's rule, and not fragile.
\item \textsc{possibly good}: supported, but narrow or fragile (few designs, a clause added to the rule after the data, or a rule chosen after the data). It appears in the main text only with its fragility stated and is never the sole proof of an innovation.
\item \textsc{inconclusive}: the numbers do not decide. An experiment whose plan fixed no threshold in advance is \textsc{inconclusive} wherever the numbers do not decide; we never invent a threshold afterwards.
\item \textsc{dismissed}: not supported (failed). Dismissed ideas get one row in the research log (Appendix~\ref{sec:log}) and are argued here only where the paper's question needs the answer.
\end{itemize}

\subsection{Statistics}
\label{sec:stats}

Ratios are geometric means over designs of the variant over the reference; below 1 favours the variant for every metric of this paper. Where a rule uses a test, the sign test counts designs that improve, the Wilcoxon test is on the paired differences (on log ratios, one-sided where stated), the Holm adjustment corrects for the family of tests the plan names, and a bootstrap interval resamples the designs (10\,000 resamples, seed 20261026; 90\% or 95\% as stated). A permutation null shuffles labels. Designs are not independent draws from a population: an interval is over this set of designs.

\subsection{M\'{\i}mir as the instrument}
\label{sec:mimir}

M\'{\i}mir~III stores, for each placed object, its provenance, its coordinates and its timing, so each experiment above can be asked of the data as well as run. What it holds today for this work: \code{structure\_mapping} rows per design, structure, run and level, with the bounding box and spread of an HDL structure's CLBs and the structure's share of the critical path (\code{crit\_share}, only with a timing run); and the run fields \code{vpr.wirelength}, \code{vpr.chan\_width}, \code{vpr.cpd}, \code{vpr.worst\_slack}, \code{vpr.logic\_clbs} and the Norns fields (\code{norns.purity}, \code{norns.reproduction}, \code{norns.cut}, \code{norns.retention}, \code{norns.critical\_matched\_delay\_share}). What is missing: near-critical statistics per HDL object (only the single worst path gives a share, so the near-critical experiment of Section~\ref{sec:e-census} cannot be recomputed from M\'{\i}mir), worst slack per structure, and power, which has no field at all.
